\section{Exámenes de 2025-26}

Preguntas del \textbf{parcial} del 31/10/2025 y del \textbf{final} del 1.\textsuperscript{er} parcial del 29/01/2026, con el sitio donde está la solución. Ninguna pregunta de examen ha quedado sin resolver: las que no son ejercicios de los boletines son resultados de los apuntes.

\begin{center}
\renewcommand{\arraystretch}{1.15}
\begin{tabular}{|l|p{6.3cm}|c|p{4.4cm}|}
\hline
\textbf{Examen} & \textbf{Pregunta} & \textbf{Pts} & \textbf{Solución} \\ \hline
Parcial, 1a & ejemplo de mismo tipo de homotopía, no homeomorfos & \multirow{2}{*}{1,5} & ejemplo tras la \cref{obs:obs_iso} \\
Parcial, 1b & contráctil $\Rightarrow$ arco-conexo, ¿recíproco? & & \textbf{T2.7} \\ \hline
Parcial, 2 & $\pi_1(X,x_0) \cong \pi_1(X,x_1)$, enunciando los auxiliares & 2 & \cref{thm:teo_cambio_base} \\ \hline
Parcial, 3 & $f \simeq$ cte $\Rightarrow f_*$ nulo, enunciando los auxiliares & 1,5 & \textbf{T3.10} \\ \hline
Parcial, 4 & existencia y unicidad de $\tilde f \colon X \to \mathbb{R}$, $X$ convexo y compacto & 3 & \cref{prop:prop_unicidad_elev} y \cref{thm:teo_existencia_elevaciones} \\ \hline
Parcial, 5 & $S^2 \setminus \{N,S\}$: recubrimientos y retracciones al ecuador (V/F) & 2 & \textbf{T4.5} \\ \hline
Final, 1a & definición de equivalencia homotópica y ejemplo que no es homeomorfismo & 0,5 & \cref{def:def_eq_homotopica} y ejemplo tras la \cref{obs:obs_iso} \\ \hline
Final, 1b & una caracterización de contráctil & 0,5 & \textbf{T2.6} \\ \hline
Final, 2 & $f_1 \simeq f_2,\ g_1 \simeq g_2 \Rightarrow g_1 \circ f_1 \simeq g_2 \circ f_2$ & 0,5 & \cref{thm:teo_comp_homo} \\ \hline
Final, 3 & definición de $\simeq_A$ y que es relación de equivalencia & 1 & \cref{thm:teo_rbe_relativa} \\ \hline
Final, 4 & producto de $\pi_1$ bien definido; inverso independiente del representante & 2 & \cref{thm:poincare} (buena definición e inverso) y \textbf{T3.1} \\ \hline
Final, 5 & enunciar $g_* = P_\sigma \circ f_*$ definiendo cada flecha; $X \simeq Y \Rightarrow \pi_1(X) \cong \pi_1(Y)$ & 1,5 & \cref{thm:aplicaciones_homotopicas} (enunciado) y \cref{prop:equiv_homotopica} \\ \hline
Final, 6 & $\deg$ bien definido e inyectivo & 2,5 & \cref{thm:teo_pi1_s1} \\ \hline
Final, 7 & toro, cilindro, recubrimiento con $\bigcap = \emptyset$, $S^2$ menos dos puntos (V/F) & 1,5 & \textbf{T4.4} \\ \hline
\end{tabular}
\end{center}
